\ifdefined\RCRSReader
  \documentclass[final,5p,times,twocolumn]{elsarticle}
\else
  \documentclass[preprint,12pt]{elsarticle}
\fi

\usepackage[T1]{fontenc}
\usepackage[utf8]{inputenc}
\ifdefined\RCRSReader\else
  \usepackage{lmodern}
\fi
\usepackage{booktabs}
\usepackage{tabularx}
\usepackage{array}
\usepackage{graphicx}
\usepackage{amsmath}
\usepackage{amssymb}
\usepackage{url}
\usepackage[hidelinks]{hyperref}

\journal{World Patent Information}
\newcommand{\RCRS}{RCRS}
\newcolumntype{Y}{>{\raggedright\arraybackslash}X}
\newcolumntype{P}[1]{>{\raggedright\arraybackslash}p{#1}}
\setlength{\emergencystretch}{2em}
\ifdefined\RCRSReader
  \usepackage{dblfloatfix}
  \usepackage{placeins}
  \newenvironment{widejournal}{\begin{table*}[!t]}{\end{table*}}
  \newcommand{\journalTableFont}{\footnotesize}
  \hypersetup{pdfsubject={WPI-style reader preview; not publisher typeset}}
\else
  \newenvironment{widejournal}{\begin{table}[!htbp]}{\end{table}}
  \newcommand{\journalTableFont}{\footnotesize}
\fi

\begin{document}

\begin{frontmatter}

\title{Retriever-Conditioned Representation Search for Cross-Domain Patent-Family Retrieval}

\author{Anonymous manuscript}

\begin{abstract}
Cross-domain prior-art retrieval must recover relevant patent families even
when the query and target use different technical language. Representation is
usually fixed before retriever evaluation, although field selection,
segmentation, packing, and family aggregation determine what a retriever can
score. We test whether representation should instead be selected conditional
on a frozen retriever. Retriever-Conditioned Representation Search (\RCRS{})
evaluates a finite grammar of executable patent-text programs without updating
retriever weights. Five lexical or dense retrievers were screened on a
family-level cross-domain benchmark. Development results varied by retriever:
program transfer was asymmetric, one search surface collapsed to a single
effective action, and fusing every primary system did not beat the strongest
single system. One 125-query selection exposure fixed two systems for an
872-query confirmation. \RCRS{} increased cross-domain (OUT) Recall@100 from 0.3311 to
0.4425 (paired difference 0.1114; 95\% bootstrap confidence interval
0.1023--0.1204), with 619 wins, 158 ties, and 95 losses. The frozen selected RCRS configuration was
then run on the complete 1,247-query benchmark. Of 5,193 strict OUT-relevant
family pairs, 796 appeared in the first 100 ranks, 332 only in ranks 101--200,
and 4,065 were absent from the Top-200 pool. Conditional representation search
therefore produced a confirmed gain under the frozen protocol, while the
post-confirmatory diagnosis placed most remaining strict OUT error outside the
exposed candidate pool.
\end{abstract}

\begin{keyword}
patent retrieval \sep patent families \sep representation search \sep
cross-domain retrieval \sep candidate exposure \sep reproducible evaluation
\end{keyword}

\end{frontmatter}

\section{Introduction}

A patent retrieval system rarely scores an invention as a single, uniform
document. Technical content may be split across family members and across
titles, abstracts, descriptions, and claims. Claims compress legal scope into
specialized language; descriptions repeat concepts at several levels of
detail. Cross-domain retrieval adds vocabulary shift between the query and the
relevant family. These properties make text construction part of the retrieval
configuration. Field order, structural labels, passage boundaries, packing,
duplicate handling, and family aggregation all affect the evidence presented
to a frozen retriever.

Family-level and embedding benchmarks now support controlled study of these
choices. DAPFAM defines domain-aware relevance at patent-family level
\cite{dapfam2026}. PatenTEB and a broader multi-task study compare patent
embeddings across retrieval, classification, and clustering tasks
\cite{patenteb2025,patentembeddings2026}. Those studies primarily compare
models or embeddings. Here the model weights stay fixed while the executable
representation changes.

AutoIndex frames representation construction as a program-search problem
\cite{autoindex2026}. Patent records require a more constrained grammar than
generic text: the program must respect document structure, long fields,
family duplication, and candidate-level aggregation. RCRS combines that
deterministic patent-family grammar with retriever-conditioned staged
evaluation and a bounded confirmation of transfer behavior. Optimization also
needs an evaluation protocol that separates iterative development from
confirmation.
Repeated access to one final set would make a representation program another
route for test-set adaptation.

Retriever-Conditioned Representation Search (\RCRS{}) addresses both points.
The experimental repository records the implementation as ArmIndex; \RCRS{}
is the publication-facing name. The method searches serializable
programs that transform patent text and aggregate family evidence. Retriever
weights, training data, and relevance judgments do not change. The protocol
separates iterative development, one Selection-125 exposure, paired Final-872
confirmation, and post-confirmatory analysis of the unchanged winner
(Fig.~\ref{fig:protocol}).

Repository receipts use A1--A7 for study records; retriever identifiers remain
in the provenance files. The article therefore uses step names and model names
in the narrative and figures; Table~\ref{tab:protocol} maps the study steps back
to their repository records.

\pagebreak
The study answers three questions. First, does one representation program
transfer across frozen retrievers, or are useful choices retriever-specific?
Second, does the selected program improve a frozen system against a frozen
static comparator on held-out confirmation? Third, after confirmation, how
much strict OUT error reflects ordering inside the retrieved pool and how much
reflects absent candidates? The claims remain within this family-level
benchmark. We do not present a new embedding model, a component-level causal
ablation, a universal baseline ranking, or a numerical comparison with
external protocols that use different corpora or relevance definitions.

\begin{figure*}[!t]
  \centering
  \includegraphics[width=\textwidth]{figures/fig1_rcrs_protocol.pdf}
  \caption{Method and study order. (A) An executable program constructs patent
  text and family-level evidence for one frozen retriever. Retriever weights
  and the scoring function are frozen; the derived candidate index is rebuilt
  deterministically for each representation program. (B) Development,
  Selection-125, Final-872 confirmation, and post-confirmatory analysis have
  separate inferential roles. System choice was fixed before Final-872.}
  \label{fig:protocol}
\end{figure*}

\section{Related work}

\subsection{Patent-family retrieval and embedding evaluation}

Patent retrieval differs from short-text web search in length, structure,
duplication, and the cost of missed prior art. BM25 remains a useful lexical
reference because it exposes exact-term evidence and has a well understood
probabilistic basis \cite{bm25review2009}. Dense retrievers add semantic
matching but impose model-specific context, pooling, and input conventions.
Multilingual and patent-specialized embedding models therefore need evaluation
at the same unit used by the application.

DAPFAM makes the family the evaluation unit and distinguishes cross-domain
relations \cite{dapfam2026}. That design avoids counting several publications
of the same invention as independent successes. PatenTEB and patent-embedding
benchmarks cover a broader set of tasks and models
\cite{patenteb2025,patentembeddings2026}. Their published scores cannot be
placed in a common ranking with the present results unless corpus version,
query construction, family normalization, relevance unit, and metric
denominators match. They position the research question but do not serve as
numerical comparators.

\subsection{Executable representations, fusion, and candidate depth}

An executable representation program records more than a prompt or a text
template. It specifies the operations that produce indexable units and the
rule that maps unit scores back to a patent family. AutoIndex provides the
closest general formulation \cite{autoindex2026}. Our grammar specializes
that idea to structured patent fields and ties every program to a frozen
retriever.

Optimization in a compound system can produce flat or unstable search
surfaces when distinct proposals compile to the same action or when a local
change is incompatible with the surrounding components
\cite{promptcoinflip2026}. The evaluation therefore records the effective
program signature and negative outcomes rather than counting proposals as
independent progress.
Rank fusion supplies a separate test of complementarity. Reciprocal rank
fusion (RRF) combines lists without calibrated scores \cite{rrf2009}, but an
extra list can still add noise. Best-single, top-two, top-three, all-primary,
and commercial-only controls were fixed before evaluation.

Candidate depth separates first-stage exposure from later ordering. Passage
retrieval can support fine-grained patent novelty analysis
\cite{finepatents2026}, but a downstream reranker or novelty model cannot act
on an absent family. The post-confirmatory oracle therefore asks only what
perfect ordering of the already frozen Top-200 could achieve. It is an
analytical bound, not a reranker result.

\section{Method}

\subsection{Retriever-conditioned objective}

Let \(a\) denote a frozen retriever, \(p\in\mathcal{P}\) an executable
representation program, \(q\) a query patent record, and \(d\) a candidate
patent-family record. Program \(p\) emits a finite set of candidate views,

\begin{equation}
 \mathcal{V}_{p}(d)=\{v_1,\ldots,v_m\}.
\end{equation}

The frozen retriever scores each view, and the program maps those scores to a
family score,

\begin{equation}
 S_{a,p}(q,d)=G_p\!\left(s_a(q,v_1),\ldots,s_a(q,v_m)\right).
\end{equation}

The development search selects

\begin{equation}
 p_a^{*}=\arg\max_{p\in\mathcal{P}}J_a(p),
\end{equation}

where \(J_a\) is the prespecified family-level objective for retriever \(a\). OUT
Recall@100 is primary; latency, cost, coverage, failures, and determinism are
guardrails. The subscript \(a\) is substantive: a program selected for one
retriever is not assumed to transfer to another.

For a query \(q\), \(R_q(q)\) denotes the query representation and \(R_p(d)\)
denotes the candidate-family views emitted by program \(p\). Retriever weights
and the scoring function are bound before evaluation. Changing \(p\) changes
the candidate views and therefore the derived index; candidate indexes are
rebuilt deterministically for each evaluated representation. The model,
scoring function, relevance judgments, and metric code remain fixed.

For a query-level relevant-family set \(\operatorname{Rel}(q)\), per-query
recall and macro query-averaged recall are
\begin{equation}
 r_q@k=\frac{|\operatorname{Rel}(q)\cap\operatorname{Top}_k(q)|}
 {|\operatorname{Rel}(q)|},\qquad
 \operatorname{Recall}@k=\frac{1}{|Q|}\sum_{q\in Q}r_q@k .
\end{equation}
The pair-level exposure diagnostic is different:
\begin{equation}
 \operatorname{Exposure}@k=\frac{\sum_q|\operatorname{Rel}(q)\cap
 \operatorname{Top}_k(q)|}{\sum_q|\operatorname{Rel}(q)|}.
\end{equation}
Macro query recall must therefore not be compared numerically with pair-level
exposure.

\subsection{Executable grammar and program identity}

The grammar contains auditable text operations. A program selects title,
abstract, and claims; orders fields; adds structural labels; creates document,
section, claim, passage, or multiview units; sets passage size, overlap, and
boundary handling; packs units within the retriever's input limit; normalizes
Unicode, whitespace, case, and punctuation; applies a duplicate policy; and
aggregates evidence by maximum score, top-three mean, top-two mean, or
view-level RRF. The program cannot change retriever parameters, training data,
the corpus, relevance labels, or metric code.

Every candidate program is serialized. Evaluation keys use its effective
action signature, so syntactically different proposals that compile to the
same behavior are counted once. This rule matters for the flat search surface
reported in Section~\ref{sec:development}.

\subsection{Finite representation-search procedure}

The search used a frozen candidate manifest rather than an unconstrained live
LLM sweep. The manifest contained 52 candidates: 40 matched candidates and 12
conditional reserves. Its proposal roles were exploit, matched-ablation,
orthogonal, and diversity, and it was frozen before measurement with seed
20260812. Each frozen proposal was compiled, mapped to an effective action
signature, and skipped if that signature had already been evaluated. For a
new signature, the procedure constructed query and candidate views, evaluated
the candidate on development data, and recorded valid, invalid, or dormant
status with cost, coverage, and failures. The predeclared decision class
(strict gain, numerical tie, or no strict gain) was then applied and the
selected configuration was frozen. No proposal mutated a retriever, corpus,
judgment, metric, or later split.

The A2 accounting was 52 frozen candidates, 44 valid measured candidates, 8
dormant reserves, and zero declared scientific failures. The separate A3
HarnessOpt run used three complete batches and 12 frozen proposals; all
proposals collapsed to one effective executable action signature, so the
search surface stopped at that flat boundary. These counts describe distinct
development records and are not additive.

\begin{widejournal}
\caption{Aggregate-safe accounting for finite representation search. The
candidate manifest was frozen before measurement; dormant reserves were not
measured.}
\label{tab:search-accounting}
\centering
\journalTableFont
\scriptsize
\setlength{\tabcolsep}{2pt}
\sloppy
\begin{tabularx}{\textwidth}{@{}P{2.8cm}P{1.0cm}P{1.9cm}P{1.5cm}P{1.1cm}Y@{}}
\toprule
Record & N & Unique sigs. & Valid & Dorm. & Decision \\
\midrule
A2 per-retriever search & 52 & not exposed & 44 & 8 & strict gain; tie; no strict gain \\
A3 HarnessOpt & 12 & 1 & -- & -- & flat-surface stop \\
\bottomrule
\end{tabularx}
\end{widejournal}

\subsection{Frozen retrievers}

All retrievers were frozen before their reported evaluations. They span a lexical
baseline, a general multilingual embedding model, a patent-specialized model,
and two recent general embedding models (Table~\ref{tab:retrievers}). Context limits
are recorded configuration properties; they do not imply equal effective
coverage of a patent record. PatEmbed uses CC BY-NC-SA 4.0, which limits
commercial deployment. The license is therefore part of the system boundary,
not a footnote to the effectiveness result.

\begin{widejournal}
\caption{Frozen retrievers. Short names are used in the text and figures;
repository retriever identifiers remain in the provenance files.}
\label{tab:retrievers}
\centering
\journalTableFont
\renewcommand{\arraystretch}{0.96}
\begin{tabularx}{\textwidth}{@{}P{1.65cm}P{2.75cm}P{2.9cm}P{2.0cm}Y@{}}
\toprule
Short name & Retriever & Type; dimension; context & License & Role in development \\
\midrule
BM25 & BM25S 0.3.10 & lexical; --; -- & -- & diagnostic lexical baseline
\cite{bm25review2009} \\
BGE-M3 & BAAI/bge-m3 & dense multilingual; 1,024; 8,192 & MIT & diagnostic
\cite{bgem3} \\
PatEmbed & datalyes/\allowbreak patembed-large & patent dense; 1,024; 512 &
CC BY-NC-SA 4.0 & advanced; research/non-commercial \\
Arctic & Snowflake Arctic-Embed-M-v2.0 & dense multilingual; 768; 8,192 &
Apache 2.0 & advanced \cite{arcticembed} \\
Qwen3 & Qwen3-Embedding-0.6B & dense multilingual; 1,024; 32,768 model-declared tokens &
Apache 2.0 & advanced \cite{qwen3embedding} \\
\bottomrule
\end{tabularx}
\end{widejournal}

\begin{widejournal}
\caption{The two systems compared on Final-872. Human-readable field,
segmentation, pooling, packing, and index settings are not included in the
safe projection; the hashes bind the frozen configurations without exposing
protected payload.}
\label{tab:frozen-systems}
\centering
\journalTableFont
\begin{tabularx}{\textwidth}{@{}P{3.0cm}P{3.0cm}Y@{}}
\toprule
System & Human-readable binding & Frozen configuration identifiers \\
\midrule
Selected RCRS & research champion; PatEmbed; research-only license &
winner binding \texttt{285e94b1...}; representation \texttt{bcbffffa...}; retrieval \texttt{e2ba107f...} \\
Static comparator & static common baseline; exact two-system comparator &
system \texttt{9a98a6d8...}; program \texttt{00f1c295...}; prompt \texttt{51a81f1b...} \\
\bottomrule
\end{tabularx}
\end{widejournal}

The ellipsized identifiers in Table~\ref{tab:frozen-systems} are display
prefixes of canonical SHA-256 values; the complete values remain in the
owner-controlled receipts. The table does not imply unreported field or
pooling settings. Final-872 is therefore an exact comparison of these two
frozen systems, not a claim against an exhaustive baseline suite.

\section{Experimental design}

\subsection{Benchmark and evidence stages}

Every experiment belongs to the same family-level cross-domain benchmark
lineage \cite{dapfam2026}. The benchmark construction is publication/source
records \(\rightarrow\) benchmark family records \(\rightarrow\) ordered text
fields \(\rightarrow\) normalization and duplicate policy \(\rightarrow\)
views/passages \(\rightarrow\) one family score. The complete retrieval pool
contains 45,336 candidate families. Table~\ref{tab:protocol} records population, exposure,
freeze state, and permitted interpretation for each study step. Development
is iterative. Selection-125 is the only selection exposure. Final-872 compares
exactly two frozen systems. The depth run and candidate-exposure diagnosis use
only the confirmed winner after Final-872.

\begin{widejournal}
\caption{Study steps, repository records, populations, and claim boundaries.
Repository codes identify experimental records, not retrievers.}
\label{tab:protocol}
\centering
\journalTableFont
\renewcommand{\arraystretch}{0.98}
\begin{tabularx}{\textwidth}{@{}P{1.2cm}P{2.3cm}P{2.6cm}P{3.4cm}Y@{}}
\toprule
Record & Study step & Population & System state and exposure & Permitted inference \\
\midrule
\mbox{A1--A3} & Development: common screen, per-retriever search, transfer and controls
& development cohorts & five retrievers; iterative search; fixed transfer and
fusion controls & describe retriever interaction, transfer, complementarity,
ties, and negative results \\
A4 & Selection-125 & 125 queries; OUT primary \(n=90\) & at most four frozen
profiles; one exposure & select one profile \\
A5 & Final-872 confirmation & 872 queries & two frozen systems; no tuning;
paired by query & estimate the confirmatory effect of \RCRS{} against the
static comparator \\
A6 & Full-benchmark depth run & 1,247 queries & confirmed winner only; depth
200; after Final-872 & report depth-wise stress test; no comparator effect \\
A7 & Candidate-exposure diagnosis & strict relation-scoped IN/OUT; OUT has 905
judged queries and 5,193 positive pairs & unchanged winner and its frozen
Top-200 pool & diagnose candidate exposure and the within-pool ordering bound \\
\bottomrule
\end{tabularx}
\end{widejournal}

\subsection{Patent-family construction and split protection}

The benchmark treats a candidate family as the unit of retrieval and
evaluation. Publication/source records are grouped into benchmark-defined
family records; ordered text fields are extracted, normalized, and subjected
to the frozen duplicate policy before views or passages are formed. View scores
are then aggregated to one family score. The safe publication projection does
not expose member identifiers, language/member-selection rules, or near-
duplicate overlap details, so those implementation settings remain an owner
gate rather than being inferred here.

The frozen split ledger is summarized below. Development comprises the
REP-DEV 150-query cohort, the HDEV-100 100-query cohort, and their combined
Train-250 250-query development record. Selection-125 contains 125 queries,
with OUT metrics on the verified (n=90) eligible subset. Final-872 contains
872 paired queries. Split membership, judgments, and hashes were frozen before
the corresponding exposure and remain protected. The safe package does not
expose query-family overlap identifiers or a zero-overlap audit; it therefore
does not claim zero cross-stage family overlap.

\begin{widejournal}
\caption{Aggregate split and exposure ledger. OUT counts are metric-population
counts, not total stage membership when otherwise stated.}
\label{tab:split-ledger}
\centering
\journalTableFont
\scriptsize
\setlength{\tabcolsep}{2pt}
\sloppy
\begin{tabularx}{\textwidth}{@{}P{1.7cm}r P{1.8cm}P{2.1cm}P{1.7cm}Y@{}}
\toprule
Split & Queries & Development record & OUT population & Exposure & Prot. \\
\midrule
Development & 150 + 100; 250 combined & REP-DEV; HDEV; Train-250 & development-specific & iterative & frozen aggregates \\
Selection-125 & 125 & -- & OUT-eligible (n=90) & one exposure & protected \\
Final-872 & 872 & -- & OUT-eligible; 17,429 pairs & held out after freeze & protected \\
\bottomrule
\end{tabularx}
\end{widejournal}

\subsection{Development, transfer, and fixed controls}

The common screen applies one representation to all five retrievers. The
per-retriever search then preserves each recorded decision class (strict
improvement, tie, or no strict improvement). The transfer study applies each
advanced-retriever program to every advanced target retriever without editing.
The \(3\times3\) matrix tests compatibility; it is not a second selection
procedure.

Five fixed controls test whether rank-list complementarity changes the result:
the strongest single system, top-two RRF, top-three RRF, all-primary fusion,
and a commercial-only profile. The all-primary and top-three controls resolve
to the same aggregate scores in the canonical artifact. None of these
development controls can replace Selection-125 or enter Final-872 after
inspection.

\subsection{Selection and confirmation}

Selection-125 evaluates the frozen research reference, BALANCED, DEEP, and FAST profiles
on 125 queries. The primary OUT subset contains 90 queries. One exposure
chooses the research profile. That profile and one static-representation
comparator are then frozen for Final-872. Confirmation requires 872/872 coverage and
paired query-level comparison; a failure would remain visible in the report.

\subsection{Metrics, uncertainty, and population accounting}

Recall@\(k\) is computed at patent-family level. OUT Recall@100 is primary.
nDCG@100 is a confirmatory secondary outcome. nDCG@10 is descriptive because
the canonical Final-872 artifact contains no confidence interval for that
endpoint. nDCG uses the benchmark's binary relevance grades. Final-872
differences are paired by query. Percentile 95\% confidence intervals use
10,000 paired bootstrap resamples of the paired query units; win, tie, and
loss counts use the same query unit. Development decision classes are strict
gain, numerical tie, and no strict gain, as recorded in the aggregate ledger.

The notation distinguishes two denominators. \(OUT_{\mathrm{eligible}}\) is
the eligible-out Final-872 cohort and contains 17,429 relevant pairs.
\(OUT_{\mathrm{strict}}\) is the relation-scoped post-confirmatory population
with 905 judged queries and 5,193 positive pairs. IN and OUT can overlap at
query level because one query may have both relation types. Final-872 and the
full-benchmark depth run therefore describe different populations; they are
not a before--after series.

\section{Development results}\label{sec:development}

\subsection{Common screen and per-retriever search}

Common-screen Recall@100 was 0.1912 for BM25, 0.2699 for BGE-M3, 0.4134 for
PatEmbed, 0.3407 for Arctic, and 0.3637 for Qwen3. The corresponding nDCG@100
values were 0.1727, 0.2314, 0.3478, 0.2845, and 0.3079. BM25 and BGE-M3
remained diagnostic; PatEmbed, Arctic, and Qwen3 entered advanced development.
Appendix Table~\ref{tab:development-ledger} gives the full common-screen and
per-retriever aggregate ledger.

The per-retriever search produced different outcomes. Arctic improved from
0.3407 to 0.3587 Recall@100 and met its strict-improvement rule. Qwen3 rose
from 0.3637 to 0.3737 but did not meet that rule. PatEmbed reached 0.4230 and
retained the recorded classification ``numerical tie to common screen.'' BM25
and BGE-M3 also tied under their decision rules, at 0.2347 and 0.2900. Program
utility therefore depended on the frozen retriever.

\subsection{Transfer, fusion, and effective search surface}

Figure~\ref{fig:transfer}A reports all nine advanced-retriever transfers. The
three programs produced a narrow Recall@100 range on PatEmbed
(0.4184--0.4193), a lower range on Arctic (0.3374--0.3413), and
0.3595--0.3626 on Qwen3. No
program source held a consistent diagonal advantage, and the target-retriever
ordering remained visible after transfer.

The diagonal cells in Fig.~\ref{fig:transfer}A are descriptive transfer cells,
not the per-retriever winner values in Appendix Table~\ref{tab:development-ledger}.
The transfer matrix and the ledger use different development cohort and
configuration summaries: the matrix applies each source program unchanged to a
target retriever, whereas the ledger reports the separate per-retriever search
record. Their diagonal values therefore need not agree, and no value is
corrected or promoted from one record to the other.

The controls were also non-monotonic (Fig.~\ref{fig:transfer}B). Best-single
and top-two RRF were nearly tied in Recall@100 (0.4184 and 0.4187), while
top-two RRF increased nDCG@100 from 0.3471 to 0.3527. All-primary fusion and
top-three RRF both returned 0.4151 Recall@100 and 0.2848 nDCG@10, below the
strongest single system on those endpoints. The commercial-only profile
returned 0.3693 Recall@100. It quantifies a licensing trade-off within this
benchmark; it is not an alternative Final winner.

The HarnessOpt procedure proposed 12 candidates in three batches, but their
effective action signatures collapsed to one. Additional proposal rounds did
not create additional executable behavior. Nominal proposal count would
therefore overstate the amount of search; the effective signature is the
reported unit.

\begin{figure*}[!b]
  \centering
  \includegraphics[width=\textwidth]{figures/fig2_transfer_controls.pdf}
  \caption{Development-only transfer and fixed controls. (A) Recall@100 for
  programs optimized on PatEmbed, Arctic, or Qwen3 and applied to each frozen
  target retriever; shading is normalized within a target column. (B) Fixed
  fusion and licensing controls. These data were not used to revise the
  Final-872 systems.}
  \label{fig:transfer}
\end{figure*}

\section{Selection and confirmatory results}

\subsection{One selection exposure}

All four Selection-125 profiles completed 125/125 queries. The research reference
obtained 0.4161 OUT Recall@100. BALANCED and DEEP tied at 0.3606; FAST
obtained 0.3083. Against either BALANCED or DEEP, the paired difference was
0.0556 (95\% CI 0.0300--0.0817; 49 wins, 20 ties, 21 losses). Against FAST,
the difference was 0.1078 (95\% CI 0.0811--0.1350; 69/9/12). The canonical
Selection-125 artifact contains no verified nDCG result, so none is inferred. The
selected RCRS configuration and static comparator were frozen after this exposure.

\subsection{Final-872 effectiveness}

Both Final-872 systems completed 872/872 queries without failure. \RCRS{} obtained
0.4425 OUT Recall@100; the static system obtained 0.3311
(Fig.~\ref{fig:final}). The paired difference was 0.1114 (95\% CI
0.1023--0.1204), with 619 wins, 158 ties, and 95 losses. nDCG@100 increased
from 0.2793 to 0.3656, a paired difference of 0.0863 (95\% CI
0.0787--0.0941). Descriptive nDCG@10 increased from 0.2337 to 0.2975,
an aggregate difference of 0.0638.

Under the frozen Final-872 protocol, conditional representation search
improved the selected system over its static comparator. The estimate applies
to this paired comparison; external retrievers and individual grammar
components were outside its design.

\begin{figure*}[!t]
  \centering
  \includegraphics[width=\textwidth]{figures/fig3_final_confirmation.pdf}
  \caption{Final-872 comparison of two frozen systems. (A) Aggregate
  effectiveness. (B) Paired query-level differences and 95\% bootstrap
  confidence intervals; nDCG@10 is descriptive. (C) Query-level Recall@100
  wins, ties, and losses.}
  \label{fig:final}
\end{figure*}

\begin{widejournal}
\caption{Final-872 effectiveness and operational checks. Confidence intervals
are paired query-level bootstrap intervals. The p99 latency result is retained
although it favors the comparator.}
\label{tab:final}
\centering
\scriptsize
\renewcommand{\arraystretch}{0.90}
\begin{tabularx}{\textwidth}{@{}P{2.4cm}r r P{1.5cm}Y@{}}
\toprule
Outcome & Static & \RCRS{} & Difference or ratio & Statistical or operational note \\
\midrule
OUT Recall@100 & 0.3311 & 0.4425 & +0.1114 &
CI 0.1023--0.1204; 619/158/95 \\
OUT nDCG@100 & 0.2793 & 0.3656 & +0.0863 &
95\% CI 0.0787--0.0941 \\
OUT nDCG@10 & 0.2337 & 0.2975 & +0.0638 &
descriptive; no canonical interval \\
Latency p50 / p95 (ms) & 315.862 / 735.171 & 237.547 / 732.218 & -- &
lower median; similar p95 \\
Latency p99 (ms) & 804.002 & 1,044.613 & -- & higher for \RCRS{} \\
Throughput & 2.665918 q/s & 3.319344 q/s & 1.245 &
direct throughput ratio (research/static) \\
Equal-cost resource metric & -- & -- & 1.456818 &
ratio computed at the same canonical cost (USD 1.456818 each); not a throughput ratio \\
Coverage / failures & 872/872; 0 & 872/872; 0 & -- &
both deterministic in the recorded run \\
\bottomrule
\end{tabularx}
\end{widejournal}

The operational results were mixed. \RCRS{} reduced median latency and had a
similar p95, but its p99 was 1,044.613 ms versus 804.002 ms for the static
system. Throughput was higher for \RCRS{} (3.319344 versus 2.665918 queries
per second), a direct ratio of 1.245. The separate equal-cost resource metric
is 1.456818 because both systems have the same canonical cost estimate of
USD 1.456818; it is not the throughput ratio. These are benchmark-scale
measurements, not evidence of patent-office deployment scalability. The worse
tail latency prevents an unqualified efficiency claim.

\section{Post-confirmatory diagnosis}

\subsection{Complete-benchmark depth profile}

After Final-872 fixed the selected RCRS configuration, the full-benchmark depth run evaluated that
system alone on 1,247 queries and 45,336 candidate families to depth 200. The run
produced 249,400 ranked rows with no failure. Overall Recall@10, @20, @50,
@100, and @200 was 0.1368, 0.2147,
0.3364, 0.4390, and 0.5468 (Fig.~\ref{fig:diagnosis}A). IN recall at those
depths was 0.1870, 0.2779, 0.4131, 0.5282, and 0.6451. Strict OUT recall was
0.0340, 0.0709, 0.1340, 0.1884, and 0.2602. nDCG@100 was 0.3625 overall,
0.4065 for IN, and 0.0706 for \ensuremath{OUT_{\mathrm{strict}}}.

The depth run contains no comparator and uses relation-scoped denominators. It describes
the frozen selected RCRS configuration at increasing candidate depth; it does not re-estimate the
Final-872 effect.

\subsection{Strict OUT exposure and Top-200 bound}

The candidate-exposure diagnosis partitions 5,193 \ensuremath{OUT_{\mathrm{strict}}}-relevant family pairs
(Fig.~\ref{fig:diagnosis}B). Of these, 796 appeared in ranks 1--100, 332 only
in ranks 101--200, and 4,065 were absent at rank 200. Among 905 judged
\ensuremath{OUT_{\mathrm{strict}}}
queries, 67 were fully retrieved by rank 200, 297 were partially retrieved,
86 had relevant evidence only in ranks 101--200, and 455 had no relevant
candidate in the Top-200 pool. Appendix Table~\ref{tab:diagnostic-ledger}
records these counts.

The analytical oracle moves each relevant family already present in the
frozen Top-200 into the first 100 positions; it adds no candidate. The
resulting Recall@100 bounds were 0.5468 overall, 0.6451 for IN, and 0.2602 for
\ensuremath{OUT_{\mathrm{strict}}}, compared with observed values of 0.4390, 0.5282, and 0.1884. The
within-pool gaps were 0.1079, 0.1169, and 0.0717
(Fig.~\ref{fig:diagnosis}C). These are ordering bounds, not achieved reranker
scores. Even perfect ordering of the exposed Top-200 would leave 4,065 strict
\ensuremath{OUT_{\mathrm{strict}}}-relevant pairs unavailable.

\section{Discussion}

\subsection{Representation is part of the retriever configuration}

The development results reject a model-independent account of preprocessing.
Arctic met its strict-improvement rule; PatEmbed and Qwen3 did not. Transfer
changed scores by target retriever and produced no consistent source-program
advantage. These outcomes fit a conditional view: the useful text units and
aggregation rule depend on the scoring behavior and input constraints of the
retriever that consumes them.

Final-872 shows that this interaction can survive a freeze. The estimate is
paired, its interval excludes zero, both systems cover the same 872 queries,
and the comparison was fixed after one selection exposure. The negative and
tied development outcomes are part of the same result. They show that
executable search can find a better configuration without guaranteeing that
every retriever, every proposal batch, or every fusion will improve.

\subsection{Candidate exposure is a separate engineering target}

The exposure decomposition narrows the next question. A reranker can reorder the
332 strict OUT-relevant pairs found only in ranks 101--200 and can change the
placement of relevant pairs already in the first 100. It cannot recover the
4,065 pairs absent from the Top-200. Candidate generation, coverage, and
cross-domain exposure therefore need separate measurement from ordering
quality.

The oracle is reported as a bound. Calling the Top-200 bound a reranking result
would imply an implemented model and an
attainable gain. Neither exists here. The bound measures only the largest
Recall@100 value permitted by the frozen candidate pool.

\subsection{Implications for patent-search evaluation}

Three reporting practices follow from the evidence. A representation program
should be versioned with its retriever, not listed as generic preprocessing.
Family normalization, duplicate handling, passage packing, and score
aggregation should appear in the executable configuration. The benchmark's
citation-derived relevance relations are evaluation labels, not legal findings
about novelty or inventive step. Finally,
relation-scoped populations should be reported with their query and pair
denominators. The same label, such as OUT, is insufficient when protocols use
different eligibility rules.

Licensing also constrains system choice. PatEmbed supports the research
comparison but its non-commercial share-alike terms may exclude a production
deployment. The commercial-only control quantifies one in-benchmark
alternative, but it does not establish the best deployable system. A
production study would need a new prespecified comparison among models with
compatible terms.

\subsection{Limitations and follow-up work}

The evidence is bounded in six ways. It comes from one family-level benchmark
with 45,336 candidate families, not a patent-office-scale deployment. The
benchmark-scale latency and throughput measurements should not be read as
production capacity evidence. Citation-derived relevance is not a legal
novelty or inventive-step determination.
Final-872 contains one frozen comparator, not an exhaustive baseline suite. The
grammar changes several components jointly, and no hash-bound component
ablation identifies a causal field, segmentation rule, or aggregation
operation. External benchmark configurations were not verified as comparable,
so cross-study scores are not ranked. IN and OUT are relation-scoped and can
overlap by query. The candidate-exposure diagnosis contains no implemented
rescue model or layer-three
reranker.

A separate preregistered study could test grammar ablations, additional licensed
baselines, an independent benchmark, and candidate-generation methods with a
new development/final split. Reusing Final-872 for those decisions would weaken the
confirmation reported here.

\ifdefined\RCRSReader\FloatBarrier\fi
\begin{figure*}[!t]
  \centering
  \includegraphics[width=\textwidth]{figures/fig4_candidate_exposure.pdf}
  \caption{Diagnosis of the unchanged selected RCRS configuration after confirmation. (A)
  Family-level recall by candidate depth. (B) Pair-level exposure of 5,193
  \(OUT_{\mathrm{strict}}\)-relevant pairs. (C) Observed Recall@100 and the analytical bound from
  perfect ordering of the frozen Top-200 pool. IN and OUT are relation-scoped
  and may overlap by query.}
  \label{fig:diagnosis}
\end{figure*}

\section{Conclusion}

Selecting an executable patent-text program for a frozen retriever increased
\(OUT_{\mathrm{eligible}}\) Recall@100 by 0.1114 over a frozen static representation on Final-872.
The gain is specific to the tested family-level protocol. After confirmation,
the unchanged selected RCRS configuration exposed only 1,128 of 5,193
\(OUT_{\mathrm{strict}}\)-relevant pairs in
its Top-200, placing most of the remaining error outside the pool available to
a reranker. The method result and the exposure limit are separate: \RCRS{}
improves representation under a freeze, while broader candidate coverage
remains an open retrieval problem.

\section*{Declaration of competing interest}

The authors declare no known competing financial interests or personal
relationships that could have influenced the work reported in this paper.

\section*{Funding}

Funding information will be inserted after author approval.

\section*{CRediT authorship contribution statement}

Author roles will be inserted after author approval and before submission.

\section*{Data and code availability}

The benchmark remains subject to its source terms, and PatEmbed remains subject
to CC BY-NC-SA 4.0. The submission package contains aggregate tables,
deterministic figure-generation code, frozen configuration identifiers, and
provenance paths for every reported number. It does not redistribute query
identifiers, patent or family identifiers, relevance judgments, per-query
outcomes, or raw rankings. The blinded repository URL, archive identifier,
and final availability statement will be added after owner review.

\section*{Declaration of generative AI and AI-assisted technologies}

During manuscript preparation, the authors used an AI-assisted drafting tool
for language and document-structure support. The authors reviewed and edited
the output, checked reported results against canonical artifacts, and retain
full responsibility for the article. This statement will be reconciled with
the journal policy in force on the submission date.

\bibliographystyle{elsarticle-num}
\bibliography{references}

\clearpage
\appendix
\ifdefined\RCRSReader\onecolumn\fi

\section{Aggregate evidence ledger}

The appendix is part of the main manuscript. No supplementary scientific file
is required. Tables~\ref{tab:development-ledger} and
\ref{tab:diagnostic-ledger} contain aggregate values only.

\begin{table}[!htbp]
\caption{Common screen (repository A1) and per-retriever search (repository
A2). Decision classes reproduce the canonical artifact and are not inferred
from displayed decimal changes.}
\label{tab:development-ledger}
\centering
\scriptsize
\setlength{\tabcolsep}{3pt}
\renewcommand{\arraystretch}{0.95}
\begin{tabularx}{\textwidth}{@{}P{1.7cm}rrr rrr Y@{}}
\toprule
& \multicolumn{3}{c}{Common screen} & \multicolumn{3}{c}{Per-retriever search} & \\
\cmidrule(lr){2-4}\cmidrule(lr){5-7}
Retriever & R@100 & nDCG@100 & nDCG@10 & R@100 & nDCG@100 & nDCG@10 & Decision class \\
\midrule
BM25 & 0.191200 & 0.172717 & 0.160011 & 0.234667 & -- & -- & diagnostic tie \\
BGE-M3 & 0.269933 & 0.231377 & 0.198497 & 0.290000 & -- & -- & diagnostic tie \\
PatEmbed & 0.413400 & 0.347812 & 0.289856 & 0.423000 & 0.357636 & 0.299444 & numerical tie \\
Arctic & 0.340667 & 0.284546 & 0.235538 & 0.358667 & 0.301868 & 0.253229 & strict gain \\
Qwen3 & 0.363733 & 0.307930 & 0.256706 & 0.373667 & 0.321262 & 0.273664 & no strict gain \\
\bottomrule
\end{tabularx}
\end{table}

\begin{table}[!htbp]
\caption{Candidate-exposure diagnosis (repository A7). The strict OUT
pair-level total is 5,193; the judged-query total is 905.}
\label{tab:diagnostic-ledger}
\centering
\journalTableFont
\renewcommand{\arraystretch}{0.95}
\begin{tabularx}{\textwidth}{@{}Y r Y r@{}}
\toprule
Pair-level class & Count & Query-level class & Count \\
\midrule
Relevant pair in Top-100 & 796 & fully retrieved by rank 200 & 67 \\
Relevant pair in ranks 101--200 & 332 & partially retrieved by rank 200 & 297 \\
Relevant pair absent at rank 200 & 4,065 & deep-only at ranks 101--200 & 86 \\
& & no relevant candidate in Top-200 & 455 \\
\bottomrule
\end{tabularx}
\end{table}

\end{document}
